\section{讨论}

门控不是把跟踪做得更好，而是选择不提交。表~\ref{tab:reconnection-main-zh} 中不限额行以下的每一个策略，接受的重连次数都更少，HOTA 的差额就是这个"不提交"的价格。把这层区别讲清楚，关系到方法应该怎么被读：HOTA 计分的是轨迹，门控计分的是决策。两者只在"门控拒绝的每一次都恰好是错的"时才一致，而数据并非如此——门控在一并拒绝错误回归的同时，也拒绝了一些正确的回归，恢复计数存在的意义就是说明它拒绝了多少。

乘积项在一个工作点有效、在另外四个无效，这是消融留下的开放问题。一个说得通的解释是：松弛天花板正是门控必须裁决边缘回归的位置，而边缘回归也正是运动一致性随间断长度衰减的地方；在严格天花板下，被接受的集合本来就小、本来就没有多少歧义，任何加权方案都没有多少可分离的空间。我们把这个解释作为解释提出，而不是作为已被检验的机制：要分清这两者，需要按提案真实难度分层，而我们的标签只记录一次已提交的决策事后是否正确。

代价应该用应用的语言来说，而不是用 HOTA 分。一个需要计人数、测停留时间、或把轨迹交给人工复核的系统，会在引发错误的那一帧之后很久还在为一次错号买单；对这样的系统，用 136 个缺口换 1,397 次失败是划算的。只按 HOTA 计分的系统不该做这个交换，表~\ref{tab:reconnection-main-zh} 已经把这一点说得很直白。

外观是明显的补充而不是竞争者。在有重识别模型和可用裁剪的场合，描述子能回答门控回答不了的问题：两个框是不是同一身衣服、同一种步态、同一件随身物品。门控提供的是另一件东西——一个不需要这些基础设施的决策版本，以及几何单独能走多远的测量。把两种信号放在同一批提案上的后续研究，可以区分"因为证据不存在而失败"和"因为评分把已有的证据称错了而失败"。我们的结果说明第二类确实存在，大约占失败量的三分之一；对第一类则什么都没说。

范围之外还有两个问题。第一，这个排序在从未被分析过的来源上是否保持。第二，同一套构造能否越过每来源十二次的上限——那是一个设计参数而不是测量结果；表~\ref{tab:caps-zh} 显示排序在各档上限下保持而量级在移动，实际部署只能根据自己的漏接容忍度来定这个上限，不能从这张表里抄一个。
